\subsection{Running the fit}

After the JAX model has been generated and the session passes readiness checks,
run parameter estimation with:

\begin{verbatim}
pfit run sessions/<session-name>
\end{verbatim}

This stage does not contact the local language model. It uses the accepted YAML
configuration, CSV files, and \texttt{generated\_script.py}. Before fitting
starts, the same deterministic readiness checks are repeated to ensure that the
generated script is current and has the required runtime interface.

Each invocation creates a new preserved run directory under:

\begin{verbatim}
sessions/<session-name>/outputs/
\end{verbatim}

The run directory contains a snapshot of the input YAML, selected data files,
generated JAX script, runtime provenance, fitted parameters, optimization
history, fitted trajectories, and summary files. These snapshots allow a
completed result to be diagnosed later even if the working session files are
changed.

The default run performs population-based global exploration followed by
gradient-based refinement. The population stage searches the bounded parameter
space in the normalized coordinates defined by the YAML configuration. The
gradient stage starts from the best population result and refines it using the
differentiable JAX loss.

For an additional refinement starting from a previous run, use:

\begin{verbatim}
pfit run sessions/<session-name> gradient-only --from-run <run-id>
\end{verbatim}

By default, the run also performs post-fit sloppiness analysis at the final
parameter estimate. This can be skipped with:

\begin{verbatim}
pfit run sessions/<session-name> --no-sloppiness
\end{verbatim}

A completed run indicates that the configured optimization procedure finished
and wrote its artifacts. It does not prove that the best fit is globally
optimal, scientifically adequate, or identifiable; those questions are examined
using the fitted trajectories, residuals, parameter locations, optimization
history, and diagnostic reports.
